\subsection{李代数的根基}

设 \(a, b\) 是李代数 \(g\) 的两个可解理想。代数 \((a + b)/b\) 同构于 \(a/(a \cap b)\)，因而是可解的；而 \(a + b\) 是由 \((a + b)/b\) 扩张 \(b\) 得到的扩张，所以也是可解的（命题 1）。由此可知，\(g\) 的一个极大可解理想包含 \(g\) 的每个可解理想，因而 \(g\) 有最大可解理想。这使我们可以作如下定义：

定义 2. 李代数的根基是它的最大可解理想。

命题 3. 李代数 \(\mathfrak{r}\) 的根基 \(\mathfrak{g}\) 是 \(\mathfrak{g}\) 的最小理想，使得 \(\mathfrak{g} / \mathfrak{r}\) 的根基为 \(\{0\}\)。

设 \(\mathfrak{a}\) 是 \(\mathfrak{g}\) 的一个理想，\(\phi\) 是将 \(\mathfrak{g}\) 映到 \(\mathfrak{g}/\mathfrak{a}\) 上的典范映射。若 \(\mathfrak{g}/\mathfrak{a}\) 的根基为零，则 \(\phi(\mathfrak{r})\) 是 \(\mathfrak{g}/\mathfrak{a}\) 的可解理想，因而为零；所以 \(\mathfrak{r} \subset \mathfrak{a}\)。另一方面，\(\overline{\phi}^{-1}(\mathfrak{r}')\) 的根基 \(\mathfrak{r}'\) 的逆像 \(\mathfrak{g}/\mathfrak{r}\) 是 \(\mathfrak{g}\) 的理想，根据命题 1 它是可解的，因而等于 \(\mathfrak{r}\)；所以 \(\mathfrak{r}' = \{0\}\)。

命题 4。设 \(\mathfrak{g}_1, \ldots, \mathfrak{g}_n\) 是李代数。各 \(\mathfrak{r}\) 的乘积的根基 \(\mathfrak{g}_i\) 等于各 \(\mathfrak{r}_i\) 的根基 \(\mathfrak{g}_i\) 的乘积。

各 \(\mathfrak{r}'\) 的乘积 \(\mathfrak{r}_i\) 是一个可解理想（命题 1），因而 \(\mathfrak{r}' \subset \mathfrak{r}\)。\(\mathfrak{r}\) 在 \(\mathfrak{g}_i\) 中的典范像是 \(\mathfrak{g}_i\) 的可解理想，因而包含于 \(\mathfrak{r}_i\)；所以 \(\mathfrak{r} \subset \mathfrak{r}'\)。

